% !TeX program = lualatex
% Standalone notebook; compile twice: lualatex 28.tex
% Research batch 39 down to 25, 27 September 2026.
% Original section preserved verbatim except for marked insertion blocks.
\documentclass[11pt,a4paper]{article}
\usepackage[margin=24mm,headheight=14pt]{geometry}
\usepackage{fontspec}
\setmainfont{Latin Modern Roman}
\setsansfont{Latin Modern Sans}
\setmonofont{Latin Modern Mono}
\usepackage{amsmath,amssymb,mathtools}
\usepackage{unicode-math}
\setmathfont{Latin Modern Math}
\usepackage{microtype}
\usepackage{enumitem,xurl,needspace}
\usepackage[dvipsnames]{xcolor}
\usepackage{hyperref}
\hypersetup{unicode=true,colorlinks=true,linkcolor=MidnightBlue,urlcolor=MidnightBlue,
 pdftitle={Research notebook: Problem 28},
 pdfauthor={Open Problems Math research notebook},bookmarksnumbered=true}
\usepackage{bookmark,titlesec,fancyhdr}
\titleformat{\section}{\Large\bfseries\raggedright}{\thesection}{0.7em}{}
\pagestyle{fancy}\fancyhf{}
\fancyhead[L]{\small\sffamily Open Problems Math\enspace|\enspace Research notebook}
\fancyhead[R]{\small\sffamily Problem 28}
\fancyfoot[L]{\small\sffamily 27 September 2026}
\fancyfoot[R]{\small\sffamily \thepage}
\setlength{\parindent}{0pt}
\setlength{\parskip}{5pt plus 1pt}
\setlength{\emergencystretch}{3em}
\setcounter{secnumdepth}{1}
\setcounter{tocdepth}{1}
\setlist[itemize]{leftmargin=3em,itemsep=5pt,topsep=4pt,parsep=0pt}
\allowdisplaybreaks
\newcommand{\N}{\mathbb{N}}
\newcommand{\Z}{\mathbb{Z}}
\newcommand{\Q}{\mathbb{Q}}
\newcommand{\R}{\mathbb{R}}
\newcommand{\C}{\mathbb{C}}
\newcommand{\F}{\mathbb{F}}
\newcommand{\A}{\mathbb{A}}
\newcommand{\PP}{\mathbb P}
\newcommand{\E}{\mathbb{E}}
\newcommand{\eps}{\varepsilon}
\newcommand{\dd}{\,\mathrm d}
\newcommand{\sref}[2]{\hyperlink{source:#1:#2}{[S#2]}}
\newcommand{\eref}[2]{\hyperlink{extra:#1:#2}{[E#2]}}
\newif\ifshowcatalogue
\showcataloguetrue
\newcommand{\cataloguescope}[1]{\ifshowcatalogue
 \paragraph{Exact scope in the supplied catalogue.}
 \begin{quote}\small #1\end{quote}\fi}
\numberwithin{equation}{section}
\begin{document}
\setcounter{section}{27}
% BEGIN PRESERVED SECTION
% ================================================================
% problemId: problem.schanuel-conjecture
% source releaseRank: 28; source releaseStatus: open
% exactTarget SHA-256: e0372e3366f5061749aad64c40e54c195891188843afac75a3ad6cdaefe5afa9
\clearpage
\section{\texorpdfstring{Schanuel's conjecture}{Schanuel's conjecture}}\label{28}
\subsection{Definitions and mathematical statement}
Numbers $z_1,\ldots,z_n$ are $\Q$-linearly independent when $\sum_jq_jz_j=0$ with $q_j\in\Q$ forces all $q_j=0$. The transcendence degree of a field extension is the maximum size of an algebraically independent subset. Schanuel's assertion is
\[
 \dim_{\Q}\operatorname{span}_{\Q}\{z_1,\ldots,z_n\}=n
 \quad\Longrightarrow\quad
 \operatorname{trdeg}_{\Q}\Q(z_1,\ldots,z_n,e^{z_1},\ldots,e^{z_n})\ge n.
\]
It is quantified over every $n\ge1$ and every such complex tuple; the $z_j$ themselves need not be algebraic.
\par\smallskip\textit{Formulation and concept references:} \sref{28}{1}.
\cataloguescope{Prove or disprove that if z\_1,...,z\_n in C are linearly independent over Q, then the field Q(z\_1,...,z\_n,e\textasciicircum{}\{z\_1\},...,e\textasciicircum{}\{z\_n\}) has transcendence degree at least n over Q.}
\subsection{Short English statement}
Do linearly independent complex numbers, together with their exponentials, always contain at least as much algebraic independence as Schanuel predicts?
% BEGIN RESEARCH 28
\subsection{Research attempt}

\paragraph{Current frontier.}
The algebraic-input case is the Lindemann--Weierstrass theorem, but the
entry quantifies arbitrary complex inputs. Waldschmidt's December 2025
manuscript proves the stronger independence of all $2n$ numbers for
almost every complex tuple, and explicitly distinguishes this from the
universal conjecture. \sref{28}{3}, Sections~1--2. Pila's recent paper
combines Schanuel and multiplicative Zilber--Pink into an equivalent joint
conjecture; it is not an unconditional proof of either. \sref{28}{2},
\eref{28}{1}. Bertolin's June 2026 revision proposes further
semi-elliptic analogues and proves special motivic cases; those hypotheses
cannot be substituted for arbitrary complex exponential tuples.
\eref{28}{2}. The catalogue's MDPI page could not be retrieved in this
run; the explicit statement supplied in the attachment and its matching
formulation in Waldschmidt were used, not an invented replacement.

\paragraph{Attack route.}
The first route is a minimal counterexample and rational change of basis:
perhaps all proper rational subspaces would force incompatible algebraic
constraints. The second is functional transcendence followed by
specialization, suggested by the abundance of generic tuples. Its missing
step is preventing transcendence degree from dropping at a specified
arithmetic point. I prove the exact minimal-counterexample constraints,
then test whether generic analytic independence actually provides that
prevention. It does not. The deductions are recorded without claiming
historical novelty for elementary predimension or specialization arguments.

\paragraph{Attempt.}
For a finite-dimensional $\Q$-subspace $W\subset\C$ with basis
$z=(z_1,\ldots,z_n)$, put
\[
 F_z=\Q(z_1,\ldots,z_n,e^{z_1},\ldots,e^{z_n}),\qquad
 d(W)=\operatorname{trdeg}_{\Q}F_z-\dim_{\Q}W.
\]
The definition of $d(W)$ must first be shown independent of its basis.

\textbf{Lemma 28.1 (basis invariance up to algebraic extension).}
If $w=Az$ for $A\in\operatorname{GL}_n(\Q)$, the compositum
$F_zF_w$ is algebraic over each of $F_z$ and $F_w$. Consequently their
transcendence degrees agree.

\emph{Proof.} Choose a positive integer $D$ clearing the denominators of
$A$. Each $w_i$ lies in $F_z$, and
\[
          (e^{w_i})^D=\prod_j(e^{z_j})^{DA_{ij}}\in F_z.
\]
Negative exponents are allowed because exponentials are nonzero. Thus
$e^{w_i}$ is algebraic over $F_z$, as an actual complex root of the
displayed equation. No arbitrary branch of logarithm is chosen. Applying
the same argument to $A^{-1}$ proves the reverse algebraicity.\hfill$\square$

\textbf{Lemma 28.2 (the form of a smallest counterexample).}
If the conjecture fails, choose a counterexample with $n$ minimal. Then
$n\ge2$,
\begin{equation}\label{r28:minimal}
                  \operatorname{trdeg}_{\Q}F_z=n-1,
\end{equation}
and for every rational hyperplane $U\subset W$, the field generated by
any basis of $U$ and its exponentials has transcendence degree $n-1$.
In a basis extending that basis of $U$, the full field is algebraic over
the hyperplane field.

\emph{Proof.} For $n=1$, a transcendental $z\ne0$ already supplies
transcendence degree at least one; for algebraic nonzero $z$, the
Lindemann theorem supplies transcendental $e^z$. This is the
one-dimensional case of the classical result discussed in
\sref{28}{3}, Section~2. Hence $n\ge2$. Extend a basis of $U$ to a basis
of $W$. By minimality, its first $n-1$ entries with their exponentials
generate a field of transcendence degree at least $n-1$. The full field
has degree at most $n-1$ by failure of the target and Lemma~28.1. Both
degrees are therefore exactly $n-1$. Since the fields are finitely
generated, equality of transcendence degrees makes the full extension
algebraic. This holds for every $U$, not only the coordinate hyperplanes.
\hfill$\square$

For $n=2$, this is quite restrictive: each nonzero rational direction
$az_1+bz_2$ must individually yield a field of transcendence degree one,
and the entire two-direction field still has degree one. But it is not
a contradiction. Many distinct generators can live in the same
transcendence-degree-one field: for instance $\Q(t)=\Q(t+a)$ for every
rational $a$ and transcendental $t$. That example does not implement the
required exponential identities; it shows why field-degree counting alone
cannot replace them. The missing step is a rigidity theorem for the
actual simultaneous additive and exponential relations in
\eqref{r28:minimal}, not a further application of the tower formula.

\medskip
\textbf{Attempt to use generic analytic independence.}
Here is a self-contained version of the generic argument, consistent
with the stronger almost-everywhere result of \sref{28}{3}.

\textbf{Lemma 28.3 (no polynomial identity of the full exponential graph).}
For every nonzero $P\in\Q[X_1,\ldots,X_n,Y_1,\ldots,Y_n]$, the entire
function
\[
                    z\longmapsto P(z,e^{z_1},\ldots,e^{z_n})
\]
is not identically zero.

\emph{Proof.} Collect monomials in $Y$ to obtain
$\sum_{a\in A}P_a(z)e^{a\cdot z}$, where $A\subset\Z_{\ge0}^n$ is finite
and each $P_a$ is nonzero. Choose a real vector $v$ so that all real
numbers $a\cdot v$ are distinct and none of the polynomials $P_a(tv)$
is identically zero. Such vectors exist: the first failures form finitely
many hyperplanes, and for each nonzero $P_a$ one nonzero homogeneous
part gives a proper algebraic exceptional set. A finite union of these
proper real algebraic sets cannot fill $\R^n$.
Restrict to $z=tv$. If the function vanished identically, divide by the
exponential with largest exponent $a\cdot v$ and let real $t\to+\infty$.
The coefficient polynomial of that term would tend to zero, since all
other polynomial terms are exponentially damped. A nonzero polynomial
cannot do so. This proves the assertion.\hfill$\square$

A nonzero holomorphic function on $\C^n$ has a Lebesgue-null zero set.
For clarity, this follows inductively from the one-variable isolated-zero
property and Fubini: expand in the last variable, choose a nonzero
holomorphic coefficient in the other variables, exclude its null zero
set, and then each remaining fiber has only a discrete zero set. Using
countably many bounded regions handles unbounded space. There are
countably many rational polynomials $P$. Thus outside a null set none
of the polynomial relations in Lemma~28.3 holds, and
\begin{equation}\label{r28:generic}
 \operatorname{trdeg}_{\Q}\Q(z_1,\ldots,z_n,e^{z_1},\ldots,e^{z_n})=2n
 \quad\hbox{for almost every }z\in\C^n.
\end{equation}
Such tuples are automatically $\Q$-linearly independent. This generic
conclusion is substantially stronger numerically than the target's
bound $n$, but weaker in its quantifier: it excludes an unspecified
null set, which can contain all the arithmetic points of interest.

\medskip
\textbf{Why specialization does not finish the proof.}
The same exponential-polynomial argument shows that $t$ and $e^t$ are
algebraically independent as functions over $\Q$. But specializing to
$t=1$ yields the field $\Q(1,e)$, of transcendence degree exactly one,
not two. Thus even a perfectly valid functional independence theorem
can lose transcendence degree at a particular algebraic specialization.
This example is compatible with Schanuel for $n=1$; it is a counterexample
to the proposed \emph{no-drop inference}, not to Schanuel.

A more concrete dependency warning comes from $(z_1,z_2)=(1,i\pi)$.
These are $\Q$-linearly independent, because one is real and the other
nonzero purely imaginary. Their exponentials are $e$ and $-1$, and
\[
 \operatorname{trdeg}_{\Q}\Q(1,i\pi,e,-1)
                  =\operatorname{trdeg}_{\Q}\Q(\pi,e).
\]
Adjoining $i$ is algebraic, which justifies the equality. The target
therefore implies algebraic independence of $e$ and $\pi$, not merely
their individual transcendence. A proposed elementary specialization
lemma that settles every tuple would have to establish this consequence
too. Numerical recognition of a few alleged polynomial relations cannot
supply a proof of either independence or its negation.

\paragraph{Stress test.}
Lemma~28.1 uses powers of actual exponentials and is insensitive to
logarithm branches; replacing it by arbitrary chosen logarithms would
require an additional $2\pi i$ term. Lemma~28.2 invokes minimality only
for strictly smaller linearly independent tuples, not for an unproved
same-dimensional case. It proves algebraicity of the missing direction
over each hyperplane field, not linearity or equality of those fields.
Neither conclusion gives a differential equation: the numbers in the
target are fixed constants, and differentiating them as if they were a
variable analytic family is illegitimate.

The generic proof explicitly permits exceptional zeros of each entire
function. Analytic continuation says a function vanishing on an open set
vanishes identically; it does not say it cannot vanish at the particular
point supplied by a hypothetical counterexample. Pila's joint-conjecture
formulation and Bertolin's motivic program articulate additional deep
structure rather than remove this obstacle by specialization.
\eref{28}{1}, \eref{28}{2}. No numerical experiment is offered as a
transcendence certificate in this entry.

\Needspace{8\baselineskip}
\paragraph{Outcome.}
\textbf{Outcome: Exploratory approach with unresolved gap.}

The attempt isolates a minimal counterexample's exact transcendence degree
and its algebraic dependence over every rational hyperplane, and proves
why generic functional independence cannot rule it out. These are useful
constraints and explicit failure tests, not a new transcendence theorem
for a previously unresolved arithmetic tuple. The unresolved task is a
rigidity or specialization argument using the exponential structure that
excludes all such minimal configurations. No complete proof or
counterexample candidate for the universal conjecture has been obtained.

% END RESEARCH 28

\subsection{Sources}
\begin{itemize}\raggedright
\item[\textnormal{[S1]}]\hypertarget{source:28:1}{} Schanuel's Conjecture and the Transcendence of Power Towers.
\url{https://www.mdpi.com/2227-7390/9/7/717}.
{\small\textit{Catalogue formulation source.}}
\item[\textnormal{[S2]}]\hypertarget{source:28:2}{} Combining the conjectures of Schanuel and Zilber-Pink.
\url{https://content.ems.press/assets/public/full-texts/serials/rlm/36/3/14299588/online/10.4171-rlm-1086.pdf}.
{\small\textit{Catalogue research/status reference; not a proof endorsement.}}
\item[\textnormal{[S3]}]\hypertarget{source:28:3}{} Schanuel Property for Elliptic and Quasi-Elliptic Functions.
\url{https://webusers.imj-prg.fr/~michel.waldschmidt/articles/pdf/SchanuelPropertyAlmostAll.pdf}.
{\small\textit{Catalogue research/status reference; not a proof endorsement.}}
% BEGIN ADDED REFERENCES 28
\item[\textnormal{[E1]}]\hypertarget{extra:28:1}{} Jonathan Pila,
\emph{Combining the conjectures of Schanuel and Zilber--Pink},
Rendiconti Lincei--Matematica e Applicazioni 36 (2025), 653--670,
DOI 10.4171/RLM/1086, abstract and Sections~1--3.
\url{https://doi.org/10.4171/RLM/1086}.
{\small\textit{Bibliographic completion of [S2]; PDF bears 2026 copyright.
Equivalence of joint conjectures, not their proof. Consulted 27 September 2026.}}
\item[\textnormal{[E2]}]\hypertarget{extra:28:2}{} Cristiana Bertolin,
\emph{A conjecture in Schanuel style for 1-motives},
arXiv:2509.08700v3, 29 June 2026.
\url{https://arxiv.org/abs/2509.08700v3}.
{\small\textit{Abstract and stated special-case scope consulted
27 September 2026; not used as an unconditional premise of a notebook lemma.}}

% END ADDED REFERENCES 28
\end{itemize}

% END PRESERVED SECTION
\end{document}
